\begin{abstract}
Este trabalho avalia o desempenho de heurísticas de Busca Local (Greedy Descent com Random Restart e Simulated Annealing) e de uma metaheurística populacional (Algoritmo Genético) na resolução do problema de sequenciamento de manutenção preventiva de ativos. O objetivo consiste em minimizar o atraso ponderado total $h(\pi)$ da operação. A avaliação foi conduzida sob um orçamento computacional fixo de avaliações da função objetivo, utilizando três instâncias de diferentes magnitudes ($n=12$, $50$ e $150$), com 30, 30 e 10 execuções independentes, respectivamente. A análise estatística foi fundamentada no teste não paramétrico de Mann-Whitney. Os resultados revelam que, embora as abordagens de busca local apresentem eficiência notável e baixo custo por iteração nas instâncias menores, o Algoritmo Genético demonstra superioridade estatística significativa na instância de grande escala, justificando seu maior custo computacional ao manter um equilíbrio otimizado entre exploração e intensificação no vasto espaço de busca.
\end{abstract}

\begin{IEEEkeywords}
Busca local, simulated annealing, algoritmo genético, sequenciamento, otimização combinatória
\end{IEEEkeywords}